% Keep the chapter heading with its opening section; preserve the locked class.
\Needspace{10\baselineskip}
\originalchapter{作业与考试的题源说明}\label{chap:assessment-map}
\originalsection{收录范围与解答来源}

本书的理论正文依据37份主讲义及1份补充讲义, 原综合练习由本书自设. 后续作业与考试章使用官方Assignments和Exams目录中的实际题面; 年度、学期、卷次和原题号在各节及题目中标明. 2015年以前印有18.440的卷不另算一门学科. 练习期末卷原题只覆盖第二次期中之后的内容, 原正式期末覆盖全课程.

官方作业目录有10份作业, 另附3页鞅说明及4道题. 作业与补充说明共134道顶层题; 其中17道Ross教材题仅登记书目定位, 不复制题面. 这些题在混合作业原件中实际有公开题面, 排除原因是外书复用限制, 不是假称无法取得题目. 收入的MIT自设题与鞅补充题为117道顶层题, 含169个正式末级子问、281项细分计算或说明输出. 一个末级子问可以要求期望和方差等多个输出, 因而这两种数量不能混用. 采访或数据采集题保留原要求; 数学解答给参数化结果和明确标注的示例, 不把未开展的采访伪造成实际观察.

考试目录有23份历年卷、2019秋季3份正式卷及1份练习期末卷, 合计27卷、199道原编号大题和582个最细子问, 包括练习卷的1道附加题. 2011春第二次期中第5题与2019秋第二次期中第5题完全相同, 后者回指前者, 同一题面和解答只收一次, 原卷定位均保留. 不同参数、不同条件或不同所求量的相近题仍分开. 一些相同小问出现在不同的大题结构中, 保留上下文并回指可复用的计算.

作业的官方目录未提供解答, 本书解答是AI补充推导. 考试有26份独立解答PDF; 2011秋第二次期中解答与题面合卷, 所以不能据独立解答文件数判断题卷遗漏. 考试解答标明依据官方答案重构的位置, 并补齐计算、证明、支持区间、边界情形与原题要求的图. 练习期末卷I.1--6、II、VI.1--6共13个子问只有参考位置或没有实质答案, 对应本书明确标注的AI补解. 其他答案中漏图、仅给计算指示以及确证公式笔误, 在对应题目附近说明.

逐资源实际页数、字节数、SHA-256、官方直链与许可边界见源码包中的两个清单: \texttt{source-manifest.json}及\texttt{course-materials-}\allowbreak\texttt{manifest.json}. 逐题及每子问的位置和解答来源见\texttt{qa/coverage-*.json}; 它们用于追踪收录, 不替代正文中的证明. 原件中的17道外书题面和受限第三方文字没有进入本书.

\originalsection{各卷原始题量}
下表按原卷出现次数计, 回指的重复题仍列入其原卷题量. 叶子指原题最细明确编号子问, 没有更细子问的整题计1; 同一子问的多个输出不另算叶子.
\begin{longtable}{@{}p{.48\linewidth}rr@{}}
\toprule 原卷 & 大题 & 末级子问\\\midrule
\endfirsthead
\toprule 原卷 & 大题 & 末级子问\\\midrule
\endhead
2011春第一次期中 & 6 & 15\\
2011秋第一次期中 & 6 & 13\\
2012秋第一次期中 & 6 & 16\\
2014春第一次期中 & 6 & 15\\
2016春第一次期中 & 6 & 18\\
2017春第一次期中 & 7 & 15\\
2018春第一次期中 & 7 & 17\\
2019春第一次期中 & 6 & 17\\
2019秋第一次期中 & 6 & 20\\\midrule
2011春第二次期中 & 6 & 16\\
2011秋第二次期中 & 5 & 17\\
2012秋第二次期中 & 6 & 18\\
2014春第二次期中 & 6 & 17\\
2016春第二次期中 & 6 & 17\\
2017春第二次期中 & 7 & 15\\
2018春第二次期中 & 7 & 18\\
2019春第二次期中 & 7 & 20\\
2019秋第二次期中 & 6 & 19\\\midrule
2011春期末 & 10 & 29\\
2012秋期末 & 10 & 30\\
2014春期末 & 10 & 28\\
2016春期末 & 10 & 31\\
2017春期末 & 10 & 30\\
2018春期末 & 10 & 33\\
2019春期末 & 10 & 33\\
2019秋期末 & 10 & 33\\
练习期末, 含附加题 & 7 & 32\\\midrule
合计 & 199 & 582\\
\bottomrule
\end{longtable}

\originalsection{补充材料与许可}
Story sheet所列分布、母函数及概念解释已融入正文相应各章, 未以英文整页附录代替中文内容. 鞅说明中的有界可选停止版本及其条件在正文第12章建立, 随附题进入作业补充节. 外部视频、文章和教材保留官方目录入口, 不逐字转录外部材料.

新增65份资源实际共436页, 其中55份、368页归档; 10份混有Ross题面的作业原件、68页整份hold, 只留来源与校验记录. 加上原讲义38份、2413页, 可分发归档为93份、2781页. 实际页数含课件渐显和空白答题空间, 不等于独立数学内容页数. MIT OCW默认CC BY-NC-SA 4.0的改编署名沿用前言说明, 可识别的第三方材料不因此自动取得相同许可.
